\section{Benchmarks}
The practical motivation for representing dependency structures as a graph, informally, 
is that computing a slice is much more performant in the graphical setting.
As we have established, backwards slicing is equivalent to computing the set of 
reachable nodes from a given selection, a procedure which runs in linear time with 
respect to the size of the reachable set. As a result, we shall see that performance 
of slicing is improved by several orders of magnitude. We shall see that 
construction of the graph does however introduce a fairly significant performance
overhead in contrast to constructing a sequential trace, however this cost vanishes when 
one computes multiple slicing queries over a single computation.

When designing interactive applications, the time taken to respond to user input is
crucial. Nielsen (1994) suggests the following guidelines:
\begin{itemize}
    \item \textbf{0.1 seconds} or less feels instantaneous to the user
    \item \textbf{1.0 seconds} is the limit before the users flow of thought feels 
    interrupted
    \item \textbf{10.0 seconds} is the limit for keeping the users attention. Anything longer
    than this and the user will shift their focus to other tasks 
\end{itemize}

We shall see that our graphical implementation pushes slicing queries into the bracket of
instantaneous user feedback. In terms of designing interactive systems, we argue that the
increased cost in constructing the dependency graph is well worth it, given that it 
facilitates such a stark increase in performance of slicing queries.

\subsection{Evaluation and Backward Slicing}
We first consider the relative performance of evaluation, and backward slicing. Table \ref{expensive-bwd} 
shows the timing data from a group of our more expensive benchmarks, measured 100s of microseconds.
\textbf{Trace-Eval} and \textbf{Trace-Bwd} measure the time taken to evaluate the program to a value
and a trace, and the time taken to run a backwards slicing query respectively. The \textbf{Graph-Eval}
and \textbf{Graph-Bwd} columns measure the time taken for the same operations but with the graphical backend.

Evaluating the program to a graph takes anywhere between 1.5 and 2 times as long as 
evaluating to a trace, but slicing queries run in less than 10 milliseconds. In
these - relatively- expensive computations, there is a concrete  performance 
benefit as soon as two or more slicing queries are made. For large systems with complex 
internal dependency structures, this is a clear benefit as it is very unlikely the
user of, for example, a large climate model will only be interested in a single slicing
criterion. Furthermore, simulations of the kind described often take a very long time
to run anyway, and so we consider the reduced performance of evaluation to be a 
worthwhile trade off.
\begin{table}
\caption{Tests: expensive, Benches: bwd}
\label{expensive-bwd}
\begin{tabular}{lrrrr}
\toprule
 & T-Eval & T-Demands & G-Eval & G-Demands \\
Test-Name &  &  &  &  \\
\midrule
slicing/convolution/edgeDetect & 346.20 & 279.00 & 1063.50 & 5.50 \\
slicing/convolution/emboss & 313.10 & 234.80 & 846.80 & 3.70 \\
slicing/convolution/gaussian & 280.40 & 230.30 & 849.70 & 7.30 \\
graphics/grouped-bar-chart & 128.80 & 107.10 & 1933.20 & 2.20 \\
graphics/line-chart & 192.10 & 158.20 & 1845.70 & 4.10 \\
graphics/stacked-bar-chart & 81.70 & 71.00 & 1255.90 & 3.70 \\
slicing/dtw/compute-dtw & 77.20 & 76.20 & 225.20 & 6.00 \\
\bottomrule
\end{tabular}
\end{table}


\subsection{Inefficient Forward Slicing}
We have already discussed that in the graphical setting, there is no theoretical need to
define a forward slicing procedure. Forward slicing is uniquely determined by 
the definition of backward slicing, since the two operators form a Galois connection.
Despite this we have formulated a more efficient forward slicing procedure. This is because
in practical cases the definition of forward slicing as the De Morgan dual of backwards 
slicing on the opposite graph can suffer from bad performance.
\small
\begin{table}
\caption{Tests: expensive, Benches: fwd}
\label{expensive-fwd}
\begin{tabular}{lrrrrr}
\toprule
 & T-Eval & T-Suffices & G-Eval & G-Suffices & Naive-Fwd \\
Test-Name &  &  &  &  &  \\
\midrule
slicing/convolution/edgeDetect & 1121.10 & 1215.10 & 3680.60 & 24.90 & 2893.40 \\
slicing/convolution/emboss & 733.40 & 640.20 & 2394.70 & 12.30 & 1469.70 \\
slicing/convolution/gaussian & 285.30 & 278.40 & 813.10 & 6.40 & 643.00 \\
graphics/grouped-bar-chart & 42.40 & 38.40 & 688.60 & 519.00 & 815.50 \\
graphics/line-chart & 62.20 & 49.70 & 533.00 & 185.10 & 678.30 \\
graphics/stacked-bar-chart & 21.90 & 20.50 & 366.70 & 277.80 & 312.40 \\
slicing/dtw/compute-dtw & 72.90 & 68.20 & 218.80 & 26.40 & 105.20 \\
\bottomrule
\end{tabular}
\end{table}


In Table \ref{expensive-fwd}, we would like to focus on the discrepancy between
\textbf{Graph-Fwd} and \textbf{Graph-FwdAsDeMorgan}. The efficient forward slicing
procedure is generally quicker than the trace based approach, but the 
version of forward slicing implemented as the De Morgan dual of backward slicing
demonstrates inferior performance. 